% Reconstruído de entradas/parte_1/Solucao_Equivalencia_Montagens_Parte_1 (1).tex



\begin{abstract}
Fixamos os tipos e as dependências de todas as construções usadas na
normalização de montagens diádicas. Definimos as representações por
átomos e por células endereçadas, suas conversões, as classes de funções,
os orçamentos e os erros em um reticulado completo. Demonstramos a
normalização por células ativas e a equivalência entre as três
formulações com os mesmos limites de folhas e altura. A ordenação finita,
a cadeia de nós e a extensão ReLU recebem construções explícitas. As
operações recursivas incluem medidas de término e casos iniciais.
\end{abstract}

\section{Base de tipos e convenções}
\label{normalizacao:sec:base}

\paragraph{Primitivos.}
Usamos os tipos padrão $\N=\{0,1,\ldots\}$, $\Z$, $\Q$ e $\R$,
as inclusões canônicas entre eles, conjuntos, subtipos, funções,
produtos, listas e conjuntos finitos. Em Lean, correspondem aos tipos
padrão e às estruturas de \code{Mathlib}. A ordem e a topologia de $\R$
são as usuais. Soma, produto, potência natural, mínimo, máximo,
continuidade, fecho, interior e cardinalidade finita têm suas definições
padrão nessa biblioteca. Os objetos próprios deste desenvolvimento são
definidos abaixo em ordem de dependência.

\paragraph{Subtipos e igualdade.}
Para $P:X\to\mathrm{Prop}$, o subtipo
$\{x:X\mid P(x)\}$ contém pares $\langle x,h\rangle$, com $h:P(x)$.
Sua projeção é $\langle x,h\rangle\mapsto x$. Dois elementos desse subtipo
são iguais quando suas projeções são iguais; as provas não constituem
dados adicionais de comparação. A igualdade de funções é extensional.

\paragraph{Conjuntos finitos.}
$\FS(X)$ é o tipo dos conjuntos finitos de elementos de $X$, sem
multiplicidades. Para $F:\FS(X)$, escrevemos $x\in F$ para pertinência e
$|F|:=\card F\in\N$. Fixamos as operações
\begin{align*}
 \operatorname{filter}(P,F)&=\{x\in F:P(x)\},\\
 \operatorname{image}(u,F)&=\{u(x):x\in F\},\\
 \operatorname{erase}(F,x)&=F\setminus\{x\}.
\end{align*}
São caracterizadas por essas equivalências de pertinência. Em particular,
filtrar produz um subconjunto, e a imagem por uma função injetiva em $F$
preserva $|F|$. Escrever $\bigcup_{x\in F}S_x$ significa o conjunto
$\{z:\exists x\in F,\ z\in S_x\}$. A finitude é mantida como dado, e
não como uma hipótese implícita sobre um conjunto arbitrário.
O conjunto subjacente a $F$ é $\{x:X:x\in F\}$.
Se $S\subseteq X$ é um conjunto finito com esse mesmo conjunto
subjacente, sua cardinalidade natural $S.\operatorname{ncard}$ é $|F|$.
Isso identifica contagens por conjuntos finitos e por \code{Finset}
quando as implementações dos postos usam representações diferentes.

\paragraph{Índices finitos.}
$\Fin(n):=\{j:\N\mid j<n\}$. Quando escrevemos $a_j$ para
$a:\Fin(n)\to X$, a desigualdade $j<n$ faz parte do índice.
Denote por $\operatorname{univ}_n:\FS(\Fin(n))$ o conjunto de todos
os elementos de $\Fin(n)$. Uma soma $\sum_{1\le j<M}$ percorre o conjunto finito
$\{j:\N:1\le j<M\}$; nela, $j$ e $j-1$ são índices válidos de
$\Fin(M)$. Se $M=1$, esse conjunto é vazio e a soma vale zero.

\paragraph{Domínio.}
Fixe $\alpha,\beta:\R$ e uma hipótese $h_K:\alpha<\beta$.
Defina separadamente o conjunto e o subtipo
\[
 K:=\{x:\R:\alpha\le x\le\beta\},\qquad
 |K|_{\rm type}:=\{x:\R\mid x\in K\},\qquad
 \iota_K:|K|_{\rm type}\to\R,\quad\iota_K(z)=z.1.
\]
Usamos $K$ também para o subtipo quando o tipo de uma função o explicita.
O intervalo aberto é $K^\circ=(\alpha,\beta)$.
Defina
\[
 \CK:=\{g:K\to\R\mid g\text{ é contínua na topologia de subtipo}\}.
\]
Para espaços topológicos $X,Y$, $C(X,Y)$ denota funções contínuas
com domínio $X$ e codomínio $Y$, acompanhadas da prova de continuidade.
Assim $\CK=C(K,\R)$ é o tipo
\code{ContinuousMap (Set.Icc alpha beta) Real}. Para $g:\CK$ e $x\in K$, a expressão $g(x)$ abrevia
$g(\langle x,h_x\rangle)$; seu valor independe da prova $h_x$.
O domínio é não vazio, compacto e igual a $\cls(K^\circ)$.

\section{Racionais, células e átomos}
\label{normalizacao:sec:objects}

% MAPA normalizacao:2.1
\begin{definicao}[racional reduzido e altura]\label{map:normalizacao:2.1}
Para $r:\Q$, sejam $p(r):\Z$ e $q(r):\N$ seu numerador e denominador
reduzidos, determinados por
\[
 q(r)>0,\qquad \gcd(|p(r)|,q(r))=1,\qquad r=p(r)/q(r).
\]
Em particular, $p(0)=0$ e $q(0)=1$. Defina
\[
 \htq:\Q\to\N,\qquad \htq(r)=\max\{|p(r)|,q(r)\},\qquad
 \qstar:=\{r:\Q\mid r\ne0\}.
\]
Aqui $|p(r)|$ é o valor absoluto natural do inteiro, isto é,
\code{Rat.num.natAbs}; $q(r)$ é \code{Rat.den}. Logo
$1\le\htq(r)$, $|p(r)|\le\htq(r)$ e $q(r)\le\htq(r)$.
\end{definicao}

% MAPA normalizacao:2.2
\begin{definicao}[célula e funções de seus extremos]\label{map:normalizacao:2.2}
O tipo de índices de células é $\Cell:=\N\times\Z$. Para $c=(k,m)$,
defina primeiro funções com valores racionais:
\begin{align*}
 L_\Q(c)&=\frac{m}{2^k},&
 R_\Q(c)&=\frac{m+1}{2^k},&
 X_\Q(c)&=\frac{2m+1}{2^{k+1}}.
\end{align*}
As operações nessas fórmulas são em $\Q$. Defina as versões reais por
inclusão: $L(c)=\uparrow L_\Q(c)$, $R(c)=\uparrow R_\Q(c)$ e
$X(c)=\uparrow X_\Q(c)$. Os três intervalos associados são
\[
 J(c)=[L(c),R(c)],\qquad O(c)=(L(c),R(c)),\qquad
 P(c)=[L(c),R(c)).
\]
$J(c)$ é o suporte fechado. $O(c)=\intr J(c)$ é seu interior.
$P(c)$ é a convenção de posse semiaberta; ela não substitui $J(c)$
nas definições de representação.
Os diádicos são
$\D:=\{\uparrow(m/2^k):(k,m)\in\Cell\}\subset\R$.
\end{definicao}

% MAPA normalizacao:2.3
\begin{lema}[geometria elementar]\label{map:normalizacao:2.3}\label{normalizacao:lem:geometry}
Para $c=(k,m)$,
\[
 R_\Q(c)-L_\Q(c)=2^{-k}>0,\qquad
 X_\Q(c)=\frac{L_\Q(c)+R_\Q(c)}2.
\]
Consequentemente, $J(c)$ é fechado e não degenerado, $O(c)$ é não
vazio e $L(c)<X(c)<R(c)$. Os três valores pertencem a $\D$.
Além disso, $J(c)=J(d)$ implica $c=d$.
\end{lema}
\begin{proof}
As duas identidades seguem da aritmética racional e $2^k>0$.
Para o centro, use o índice $(k+1,2m+1)$ na definição de $\D$.
Se os suportes são iguais, seus mínimos e máximos são iguais.
As larguras dão $2^{-k}=2^{-k'}$, logo $k=k'$, pela estrita
monotonicidade das potências de dois. A igualdade dos extremos
esquerdos dá então $m=m'$.
\end{proof}

% MAPA normalizacao:2.4
\begin{definicao}[átomo e suas projeções]\label{map:normalizacao:2.4}\label{normalizacao:def:atom}
$\Atoms$ é a estrutura com campos
\[
 k:\N,\quad m:\Z,\quad y:\Q,\quad s:\Q,\quad h_s:s\ne0.
\]
A igualdade de átomos é a igualdade dos quatro campos de dados
$(k,m,y,s)$. Para $a:\Atoms$, defina
\begin{align*}
 C(a)&=(a.k,a.m):\Cell, & J_a&=J(C(a)), & O_a&=O(C(a)),\\
 L_a&=L(C(a)),&R_a&=R(C(a)),&X_a&=X(C(a)),\\
 b(a)&=a.s:\Q,& u(a)&=a.y-a.s\,X_\Q(C(a)):\Q.
\end{align*}
\begin{align*}
 \ell_a(x)&=\uparrow a.y+\uparrow a.s\,(x-X_a)\quad(x:\R),\\
 h(a)&=\max\{2,a.k+1,|a.m|+1,\htq(a.y),\htq(a.s)\}:\N.
\end{align*}
O átomo de valor padrão é
$a_\dagger=(0,0,0,1,\,1\ne0)$.
\end{definicao}

% MAPA normalizacao:2.5
\begin{lema}[formas da reta]\label{map:normalizacao:2.5}\label{normalizacao:lem:line}
Para todo $a:\Atoms$ e $x:\R$,
\[
 \ell_a(x)=\uparrow u(a)+\uparrow b(a)\,x,\qquad
 \ell_a(X_a)=\uparrow a.y.
\]
A função $\ell_a:\R\to\R$ é contínua e tem inclinação não nula.
Se $C(a)=C(a')$, $a.y=a'.y$ e $a.s=a'.s$, então
$a=a'$, $J_a=J_{a'}$ e $\ell_a=\ell_{a'}$ como funções globais.
\end{lema}
\begin{proof}
Expanda os produtos, usando que a inclusão $\Q\to\R$ preserva soma
e produto. A continuidade é a de uma função afim. A inclusão é
injetiva, logo $h_s$ também garante a inclinação real não nula.
As igualdades finais seguem da igualdade dos campos da estrutura.
\end{proof}

\section{Predicados de representação e normalização}
\label{normalizacao:sec:normalization}

% MAPA normalizacao:3.1
\begin{definicao}[campos dos contratos]\label{map:normalizacao:3.1}\label{normalizacao:def:contracts}
Para $\Lambda:\FS(\Atoms)$, $g:\CK$ e $G:\R\to\R$, defina
\begin{align*}
 \Sep(\Lambda)&:\Longleftrightarrow
  \forall a,b\in\Lambda,\ a\ne b\Rightarrow O_a\cap O_b=\varnothing,\\
 \Cov(K,\Lambda)&:\Longleftrightarrow
  \forall x\in K,\ \exists a\in\Lambda,\ x\in J_a,\\
 \Match_K(g,\Lambda)&:\Longleftrightarrow
  \forall a\in\Lambda\ \forall z:K,\
       \iota_K(z)\in J_a\Rightarrow g(z)=\ell_a(\iota_K(z)),\\
 \Match^G_K(G,\Lambda)&:\Longleftrightarrow
  \forall a\in\Lambda\ \forall x\in J_a\cap K,\ G(x)=\ell_a(x),\\
 \Agree(\Lambda)&:\Longleftrightarrow
  \forall a,b\in\Lambda\ \forall x\in J_a\cap J_b,\
       \ell_a(x)=\ell_b(x).
\end{align*}
Defina os predicados, todos com valores em $\mathrm{Prop}$,
\begin{align*}
 \Loc_K(g,\Lambda)&:=\Sep(\Lambda)\land\Cov(K,\Lambda)
                         \land\Match_K(g,\Lambda),\\
 \Str_K(g,\Lambda)&:=\Loc_K(g,\Lambda)\land\Agree(\Lambda),\\
 \Loc^G_K(G,\Lambda)&:=\Sep(\Lambda)\land\Cov(K,\Lambda)
                         \land\Match^G_K(G,\Lambda),\\
 \Str^G_K(G,\Lambda)&:=\Loc^G_K(G,\Lambda)\land\Agree(\Lambda).
\end{align*}
Finalmente, como subconjuntos do mesmo tipo $\CK$,
\[
 \Aloc(K):=\{g:\CK:\exists\Lambda,\ \Loc_K(g,\Lambda)\},\qquad
 \Astr(K):=\{g:\CK:\exists\Lambda,\ \Str_K(g,\Lambda)\}.
\]
\end{definicao}

A igualdade das retas em $J_a\cap J_b\cap K$ é consequência de
$\Match$: ambas têm o valor $g(x)$. A cláusula $\Agree$ exige
igualdade também nos encontros fora de $K$.

% MAPA normalizacao:3.2
\begin{definicao}[atividade e filtro]\label{map:normalizacao:3.2}\label{normalizacao:def:norm}
Defina $\Act_K:\Atoms\to\mathrm{Prop}$ e
$\Norm_K:\FS(\Atoms)\to\FS(\Atoms)$ por
\[
 \Act_K(a):\Longleftrightarrow O_a\cap K^\circ\ne\varnothing,
 \qquad \Norm_K(\Lambda):=\operatorname{filter}(\Act_K,\Lambda).
\]
Portanto,
$a\in\Norm_K(\Lambda)\Longleftrightarrow a\in\Lambda\land\Act_K(a)$.
Atividade significa interseção não vazia dos interiores.
\end{definicao}

% MAPA normalizacao:3.3
\begin{lema}[critério exato de atividade]\label{map:normalizacao:3.3}\label{normalizacao:lem:active}
Para todo $a:\Atoms$,
\[
 \Act_K(a)\Longleftrightarrow L_a<\beta\ \land\ \alpha<R_a.
\]
\end{lema}
\begin{proof}
Um ponto da interseção implica as duas desigualdades. Reciprocamente,
ponha $v=\max\{L_a,\alpha\}$ e $w=\min\{R_a,\beta\}$. De
$L_a<R_a$, $\alpha<\beta$ e das duas desigualdades segue $v<w$.
O ponto $(v+w)/2$ pertence a ambos os intervalos abertos.
\end{proof}

\paragraph{Decidibilidade.}
Para $\alpha,\beta\in\R$, o filtro é definido em Lean usando
\code{Classical.propDecidable}; isso é uma escolha lógica de
instâncias, não uma rotina efetiva de comparação de reais.
Se os extremos são racionais, o lema~\ref{normalizacao:lem:active} fornece um filtro
executável por comparações racionais exatas.

% MAPA normalizacao:3.4
\begin{lema}[intervalos separados]\label{map:normalizacao:3.4}\label{normalizacao:lem:separated}
Se $u<v$, $p<q$ e $(u,v)\cap(p,q)=\varnothing$, então
$v\le p$ ou $q\le u$. Se $[u,v]\cap[p,q]$ é não vazio, ele é
uma única extremidade comum.
\end{lema}
\begin{proof}
A negação das duas desigualdades dá
$\max\{u,p\}<\min\{v,q\}$. Seu ponto médio contradiz a
disjunção. No caso $v\le p$, uma interseção implica $v=p$ e
ela é $\{v\}$. O outro caso é simétrico.
\end{proof}

% MAPA normalizacao:3.5
\begin{teorema}[normalização]\label{map:normalizacao:3.5}\label{normalizacao:thm:norm}
Se $\Loc_K(g,\Lambda)$, então, para $\Lambda'=\Norm_K(\Lambda)$,
\begin{enumerate}
\item $\Lambda'\subseteq\Lambda$ e $|\Lambda'|\le|\Lambda|$;
\item $\Cov(K,\Lambda')$ e $\Lambda'\ne\varnothing$;
\item $J_a\cap J_b\subseteq K^\circ$ para $a\ne b$ em $\Lambda'$;
\item $\Str_K(g,\Lambda')$;
\item $\Norm_K(\Lambda')=\Lambda'$.
\end{enumerate}
As mesmas conclusões, com os predicados globais, valem sob
$\Loc^G_K(G,\Lambda)$; nesse caso a função $G$ inteira é preservada.
\end{teorema}
\begin{proof}
A inclusão e a cota de cardinalidade são propriedades de filtro.
Se $x\in K^\circ$ e $x\in J_a$ com $a\in\Lambda$, então
$L_a\le x<\beta$ e $\alpha<x\le R_a$. O
lema~\ref{normalizacao:lem:active} implica $\Act_K(a)$. Logo
$K^\circ\subseteq U:=\bigcup_{a\in\Lambda'}J_a$.
$U$ é uma união finita de fechados e, portanto, é fechado.
Assim $K=\cls(K^\circ)\subseteq U$. Como $\alpha\in K$,
a família $\Lambda'$ não é vazia.

Para $a\ne b$ retidos, o lema~\ref{normalizacao:lem:separated} mostra que
uma interseção não vazia é $\{z\}$, com
$R_a=z=L_b$ após eventual troca dos nomes. Escolha pontos ativos
$u\in O_a\cap K^\circ$ e $v\in O_b\cap K^\circ$. Então
$\alpha<u<z<v<\beta$, provando a terceira conclusão.

Disjunção dos interiores e coincidência com $g$ são herdadas da
família original. Em cada interseção entre átomos distintos retidos,
a terceira conclusão permite escrever
$\ell_a(z)=g(z)=\ell_b(z)$. Para átomos iguais, a igualdade é
reflexiva. Obtemos $\Agree(\Lambda')$ e, portanto, $\Str$.
A idempotência segue da equivalência de pertinência do filtro:
$P\land P\Longleftrightarrow P$. O argumento global usa
$G(z)$ no lugar de $g(z)$ e não altera $G$ em nenhum ponto.
\end{proof}

% MAPA normalizacao:3.6
\begin{corolario}[equivalência local e consolidada]\label{map:normalizacao:3.6}\label{normalizacao:cor:classes}
$\Aloc(K)=\Astr(K)$ como subconjuntos de $\CK$.
\end{corolario}
\begin{proof}
Esquecer $\Agree$ dá uma inclusão. Aplicar o
teorema~\ref{normalizacao:thm:norm} a uma testemunha local dá a outra.
\end{proof}

\section{Enumerações fixadas por fórmulas}
\label{normalizacao:sec:enumeration}

As enumerações usadas nesta seção são determinadas pelas fórmulas do
pacote examinado. Uma sobrejeção abstrata não substitui esses dados.

% MAPA normalizacao:4.1
\begin{definicao}[pareamento e descompactação]\label{map:normalizacao:4.1}
Defina $T:\N\to\N$ por recursão
\[
 T(0)=0,\qquad T(n+1)=T(n)+n+1.
\]
Assim $T(n)=n(n+1)/2$. Defina
$\pi(u,v)=T(u+v)+v$. Para $n:\N$, o conjunto finito
$\{j\in\{0,\ldots,n\}:T(j)\le n\}$ contém zero. Ponha
\begin{align*}
 w(n)&=\max\{j\le n:T(j)\le n\},\\
 v(n)&=n-T(w(n)),& u(n)&=w(n)-v(n),&
 \pi^{-1}(n)&=(u(n),v(n)).
\end{align*}
As subtrações são naturais; as cotas provadas abaixo garantem que não
há truncamento na região usada. Finalmente, defina
\begin{align*}
 \operatorname{pack}(k,\mu,r,t)&=\pi(\pi(k,\mu),\pi(r,t)),\\
 \operatorname{unpack}(n)&=(k,\mu,r,t),
 \quad\text{onde }(u,v)=\pi^{-1}(n),\ (k,\mu)=\pi^{-1}(u),\
 (r,t)=\pi^{-1}(v).
\end{align*}
Esses mapas têm tipos $\N^4\to\N$ e $\N\to\N^4$.
\end{definicao}

% MAPA normalizacao:4.2
\begin{lema}[inversas do pareamento]\label{map:normalizacao:4.2}\label{normalizacao:lem:pair}
$T(w(n))\le n<T(w(n)+1)$, $0\le v(n)\le w(n)$ e
\[
 \pi^{-1}(\pi(u,v))=(u,v),\quad
 \pi(\pi^{-1}(n))=n.
\]
Consequentemente, $\operatorname{pack}$ e $\operatorname{unpack}$
são inversas. A expressão de $w$ usada no código é
\[
 w(n)=\left\lfloor
 \frac{\lfloor\sqrt{8n+1}\rfloor-1}{2}\right\rfloor,
\]
com raiz quadrada inteira \code{Nat.sqrt} e divisão natural.
\end{lema}
\begin{proof}
A recorrência dá $T(j+1)-T(j)=j+1>0$ e $T(j)\ge j$.
A definição de máximo fornece a primeira desigualdade. Se
$T(w+1)\le n$, então $w+1\le n$, contradizendo a maximalidade;
logo vale a segunda. Portanto $0\le n-T(w)<w+1$ e $v\le w$.
Como $u+v=w$, temos $\pi(u,v)=T(w)+v=n$.
Para um par dado, $T(u+v)\le\pi(u,v)<T(u+v+1)$ determina
unicamente a diagonal e, nela, $v$. Isso prova as duas identidades.
As inversas dos mapas em quatro coordenadas seguem por três aplicações.
Por fim, $8T(w)+1=(2w+1)^2$; as desigualdades da diagonal dão
$2w+1\le\sqrt{8n+1}<2w+3$, que implica a fórmula inteira.
\end{proof}

% MAPA normalizacao:4.3
\begin{definicao}[enumeração dos inteiros]\label{map:normalizacao:4.3}
Defina $Z:\N\to\Z$ por casos de paridade e $Z^{-1}:\Z\to\N$ por
\[
 Z(2j)=j,\qquad Z(2j+1)=-(j+1),\qquad
 Z^{-1}(m)=\begin{cases}2m,&m\ge0,\\-2m-1,&m<0.\end{cases}
\]
Os valores da última expressão são naturais. Substituir os dois casos
mostra $Z(Z^{-1}(m))=m$ e $Z^{-1}(Z(n))=n$.
\end{definicao}

% MAPA normalizacao:4.4
\begin{definicao}[ordens e postos dos racionais]\label{map:normalizacao:4.4}\label{normalizacao:def:rank}
Para $H:\N$, defina o conjunto finito
\[
 \mathcal Q_H=
 \left\{\frac{p}{q}:p\in\Z,\ |p|\le H,
                         \ q\in\N,\ 1\le q\le H\right\}.
\]
Como conjunto de valores racionais, ele contém cada valor uma única
vez. É exatamente $\{r:\Q:\htq(r)\le H\}$: reduzir uma fração não
amplia numerador ou denominador, e a representação reduzida dá a
inclusão inversa.

A ordem $\prec_2$, também usada para $\prec_3$, é a ordem
lexicográfica da tripla
\[
 \bigl(\htq(r),q(r),p(r)\bigr):\N\times\N\times\Z.
\]
Na ordem lexicográfica, compara-se a primeira coordenada diferente.
Para a implementação consolidada, defina a chave natural
\[
 \kappa(r)=4\htq(r)^3+q(r)(2\htq(r)+1)
                  +(p(r)+\htq(r))_{\N},
 \qquad r\prec_{\rm c}s:\Longleftrightarrow\kappa(r)<\kappa(s).
\]
O símbolo $(z)_{\N}$ é \code{Int.toNat}; neste uso,
$p(r)+\htq(r)\ge0$.
Para $i\in\{2,3,\mathrm c\}$, defina
\begin{align*}
 \pred_i(r):\FS(\Q)&=\{v\in\mathcal Q_{\htq(r)}:v\prec_i r\},
 &\rankq_i(r)&=|\pred_i(r)|,\\
 \pred_i^*(r):\FS(\qstar)&=\{v:\qstar:v.1\in\mathcal Q_{\htq(r.1)},\ v.1\prec_i r.1\},
 &\rankq_i^*(r)&=|\pred_i^*(r)|.
\end{align*}
Os tipos dos postos são $\Q\to\N$ e $\qstar\to\N$.
\end{definicao}

% MAPA normalizacao:4.5
\begin{lema}[postos e enumerações inversas]\label{map:normalizacao:4.5}\label{normalizacao:lem:rank}
As três ordens da definição~\ref{normalizacao:def:rank} coincidem e cada conjunto
de predecessores é finito e contém todos os predecessores na ordem
respectiva. Cada $\rankq_i$ e $\rankq_i^*$ é bijetiva.
Definimos, portanto, univocamente
\[
 Q_i:\N\to\Q,\quad\rankq_i(Q_i(n))=n,
 \qquad Q_i^*:\N\to\qstar,\quad\rankq_i^*(Q_i^*(n))=n.
\]
Valem também $Q_i(\rankq_i(r))=r$ e
$Q_i^*(\rankq_i^*(s))=s$.
\end{lema}
\begin{proof}
Para altura $H\ge1$, a parte da chave $\kappa$ posterior a $4H^3$
está entre zero e $2H^2+3H$, que é menor que
$4(H+1)^3-4H^3$. Logo alturas diferentes são comparadas primeiro
pela altura. A altura fixa, aumentar o denominador em um aumenta
a chave em pelo menos $(2H+1)-2H=1$, e a denominador fixo o
numerador decide a ordem. Assim $\prec_{\rm c}$ é precisamente a
ordem lexicográfica. Qualquer predecessor de $r$ tem altura no
máximo $\htq(r)$, justificando os filtros finitos.

Em qualquer ordem linear infinita com conjuntos finitos de
predecessores, o posto dado por sua cardinalidade é estritamente
crescente: se $a\prec b$, os predecessores de $a$ estão entre os de
$b$, e $a$ pertence somente ao segundo conjunto. Logo o posto é
injetivo. Sua imagem é infinita, portanto não limitada em $\N$.
Dado $n$, escolha $b$ com posto pelo menos $n$. O conjunto finito
ordenado $\pred(b)\cup\{b\}$ tem um elemento na posição $n$,
contando a partir de zero. Todos os predecessores desse elemento
pertencem a esse conjunto, e são exatamente os $n$ elementos
anteriores; seu posto é $n$. Isso prova sobrejetividade.
A ordem de $\Q$ é infinita pela inclusão de $\Z$; a ordem restrita a
$\qstar$ é infinita pela inclusão dos inteiros positivos.
As identidades das inversas decorrem de existência e unicidade.
\end{proof}

% MAPA normalizacao:4.6
\begin{definicao}[codificadores concretos de átomos]\label{map:normalizacao:4.6}\label{normalizacao:def:codec}
Para cada $i\in\{2,3,\mathrm c\}$, defina
\begin{align*}
 \enc_i(a)&=\operatorname{pack}
    \bigl(a.k,Z^{-1}(a.m),\rankq_i(a.y),
                     \rankq_i^*(\langle a.s,h_s\rangle)\bigr),\\
 \dec_i(n)&=
 \bigl(k,Z(\mu),Q_i(r),(Q_i^*(t)).1,(Q_i^*(t)).2\bigr),
 \qquad (k,\mu,r,t)=\operatorname{unpack}(n).
\end{align*}
Assim $\enc_i:\Atoms\to\N$ e $\dec_i:\N\to\Atoms$.
Defina ainda
\[
 C_i(n)=C(\dec_i(n)),\qquad
 U_i(n)=u(\dec_i(n)),\qquad B_i(n)=b(\dec_i(n)),
\]
\[
 \lambda_i(n,x)=\uparrow U_i(n)+\uparrow B_i(n)\,x
 \quad(n:\N,\ x:\R).
\]
\end{definicao}

% MAPA normalizacao:4.7
\begin{lema}[contrato completo de endereçamento]\label{map:normalizacao:4.7}\label{normalizacao:lem:codec}
Para cada $i$,
\[
 \dec_i(\enc_i(a))=a,\qquad \enc_i(\dec_i(n))=n,
 \qquad\lambda_i(n,x)=\ell_{\dec_i(n)}(x).
\]
Para $i,j\in\{2,3,\mathrm c\}$, o reendereçamento é a função
$\tau_{i\to j}:=\enc_j\circ\dec_i:\N\to\N$. Ela satisfaz
\[
 \dec_j(\tau_{i\to j}(n))=\dec_i(n),\qquad
 C_j(\tau_{i\to j}(n))=C_i(n),\qquad
 \lambda_j(\tau_{i\to j}(n),x)=\lambda_i(n,x).
\]
\end{lema}
\begin{proof}
As duas primeiras identidades resultam das inversas de
$\operatorname{pack}$, $Z$ e dos postos racionais, coordenada a
coordenada. A identidade de retas é o lema~\ref{normalizacao:lem:line}.
Substituir $a=\dec_i(n)$ na primeira identidade prova todas as
propriedades de $\tau$.
\end{proof}

\section{Representações por células e mapas entre os tipos de função}
\label{normalizacao:sec:bridge}

% MAPA normalizacao:5.1
\begin{definicao}[testemunha por células endereçadas]\label{map:normalizacao:5.1}\label{normalizacao:def:cellrep}
Para $F:\FS(\Cell)$ e $e:\Cell\to\N$, defina
$\operatorname{CellRep}_{i,K}(g,F,e)$, com $g:\CK$, pelas quatro
cláusulas
\begin{align*}
 &\forall c,d\in F,\ c\ne d\Rightarrow O(c)\cap O(d)=\varnothing;\\
 &\forall x\in K,\ \exists c\in F,\ x\in J(c);\\
 &\forall c\in F,\ C_i(e(c))=c;\\
 &\forall c\in F\ \forall z:K,\ \iota_K(z)\in J(c)
       \Rightarrow g(z)=\lambda_i(e(c),\iota_K(z)).
\end{align*}
O predicado $\operatorname{CellRep}^G_{i,K}(G,F,e)$ tem as mesmas
três primeiras cláusulas e substitui a última por
\[
 \forall c\in F\ \forall x\in J(c)\cap K,\quad
 G(x)=\lambda_i(e(c),x),\qquad G:\R\to\R.
\]
\end{definicao}

% MAPA normalizacao:5.2
\begin{definicao}[continuidade no domínio e restrição]\label{map:normalizacao:5.2}\label{normalizacao:def:restriction}
Para $G:\R\to\R$, $\operatorname{ContOn}_K(G)$ significa
\[
 \forall x\in K\ \forall\varepsilon>0\ \exists\delta>0\
 \forall y\in K,\quad |y-x|<\delta\Rightarrow
                         |G(y)-G(x)|<\varepsilon.
\]
Sob uma prova $h:\operatorname{ContOn}_K(G)$,
$\Res_K(G,h):\CK$ é a função $z\mapsto G(\iota_K(z))$ com a
continuidade garantida por $h$.
Para $\mathcal S\subseteq(\R\to\R)$, defina
\[
 \ImageRes_K(\mathcal S)=
 \{g:\CK:\exists G\in\mathcal S,\
             \forall z:K,\ g(z)=G(\iota_K(z))\}.
\]
Essa imagem é um subconjunto de $\CK$ e não depende da escolha
de uma prova de continuidade.
\end{definicao}

% MAPA normalizacao:5.3
\begin{definicao}[as três classes tipadas]\label{map:normalizacao:5.3}\label{normalizacao:def:threeclasses}
Fixe exatamente
\begin{align*}
 \mathcal B_2(K)&:=
 \{g:\CK:\exists F:\FS(\Cell)\ \exists e:\Cell\to\N,\
                  \operatorname{CellRep}_{2,K}(g,F,e)\},\\
 \mathcal B_3(K)&:=
 \{G:\R\to\R:\operatorname{ContOn}_K(G)\ \land\\
 &\hspace{32mm}\exists F:\FS(\Cell)\ \exists e:\Cell\to\N,\
                   \operatorname{CellRep}^G_{3,K}(G,F,e)\},\\
 \mathcal B_{\rm con}(K)&:=
 \{G:\R\to\R:\operatorname{ContOn}_K(G)\ \land\
                  \exists\Lambda:\FS(\Atoms),\ \Str^G_K(G,\Lambda)\}.
\end{align*}
\end{definicao}

% MAPA normalizacao:5.4
\begin{lema}[conversões explícitas das testemunhas]\label{map:normalizacao:5.4}\label{normalizacao:lem:witnesses}
As testemunhas por células e por átomos se convertem com a mesma
cardinalidade e com as mesmas peças afins, antes da normalização.
As funções de conversão são as seguintes.
\end{lema}
\begin{proof}
\emph{Células para átomos.} Dada $(F,e)$ válida para a fase $i$,
defina
\[
 T_i(F,e):=\operatorname{image}
             (c\mapsto\dec_i(e(c)),F):\FS(\Atoms).
\]
A compatibilidade $C_i(e(c))=c$ mostra que esse mapa é injetivo em
$F$: a igualdade dos átomos implica a igualdade das células.
Logo $|T_i(F,e)|=|F|$. As cláusulas de disjunção, cobertura e
coincidência com $g$ são transportadas pelo lema~\ref{normalizacao:lem:codec}.

\emph{Átomos para células.} Suponha $\Sep(\Lambda)$. O mapa
$C$ é injetivo em $\Lambda$. De fato, se $a\ne b$ tivessem a
mesma célula, então $O_a=O_b\ne\varnothing$, contradizendo
$\Sep$. Defina
\[
 I(\Lambda):=\operatorname{image}(C,\Lambda):\FS(\Cell).
\]
Para todo $c:\Cell$, defina a função total
\[
 \sigma_\Lambda(c)=
 \begin{cases}
 \text{o único }a\in\Lambda\text{ tal que }C(a)=c,
                                      &c\in I(\Lambda),\\
 a_\dagger,                            &c\notin I(\Lambda).
 \end{cases}
\]
A unicidade foi provada, e a existência é a equivalência de
pertinência da imagem. Em Lean, essa é uma descrição única sob
\code{Classical}; nenhum átomo novo é introduzido na família.
Defina a função de endereçamento total
\[
 e_{i,\Lambda}:\Cell\to\N,\qquad
 e_{i,\Lambda}(c):=\enc_i(\sigma_\Lambda(c)).
\]
Para $c\in I(\Lambda)$, a decodificação recupera precisamente o
único átomo original da célula $c$. Consequentemente,
$(I(\Lambda),e_{i,\Lambda})$ satisfaz as cláusulas celulares,
$|I(\Lambda)|=|\Lambda|$ e
\[
 T_i(I(\Lambda),e_{i,\Lambda})=\Lambda.
\]
Na direção oposta, $I(T_i(F,e))=F$, e o endereçamento reconstruído
coincide com $e$ em $F$, pela identidade $\enc_i\circ\dec_i=\mathrm{id}$.
As mesmas construções e provas usam $\Match^G$ para funções globais.
\end{proof}

% MAPA normalizacao:5.5
\begin{definicao}[extensão contínua por projeção]\label{map:normalizacao:5.5}\label{normalizacao:def:clamp}
Defina o mapa $p_K:\R\to K$ e a extensão $E_0:\CK\to C(\R,\R)$ por
\[
 p_K(x)=\langle\min\{\beta,\max\{\alpha,x\}\},h_x\rangle,
 \qquad E_0(g)(x)=g(p_K(x)).
\]
A prova $h_x$ é dada pelas desigualdades
$\alpha\le\min\{\beta,\max\{\alpha,x\}\}\le\beta$.
Mínimo e máximo são contínuos, logo $p_K$ é contínuo como mapa
para o subtipo e $E_0(g)$ é contínua por composição.
Se $x\in K$, então $\iota_K(p_K(x))=x$; portanto
$\Res_K(E_0(g))=g$.
\end{definicao}

% MAPA normalizacao:5.6
\begin{teorema}[equivalência extensional tipada]\label{map:normalizacao:5.6}\label{normalizacao:thm:bridge}
Como subconjuntos de $\CK$,
\begin{align*}
 \mathcal B_2(K)&=\Aloc(K)=\Astr(K),\\
 \ImageRes_K(\mathcal B_3(K))&=\Aloc(K),\\
 \ImageRes_K(\mathcal B_{\rm con}(K))&=\Astr(K).
\end{align*}
Como subconjuntos de $\R\to\R$,
$\mathcal B_3(K)=\mathcal B_{\rm con}(K)$.
\end{teorema}
\begin{proof}
O lema~\ref{normalizacao:lem:witnesses} identifica $\mathcal B_2(K)$ com
$\Aloc(K)$ e traduz as testemunhas globais da Fase 3 para
$\Loc^G$. A restrição de cada função global com continuidade em
$K$ produz um elemento de $\CK$ com a mesma representação.
Reciprocamente, $E_0(g)$ é uma função global contínua com a mesma
coincidência com as peças em $K$. Isso prova as duas igualdades
de imagens de restrição. O corolário~\ref{normalizacao:cor:classes} identifica
as classes local e consolidada no tipo $\CK$.

Para a igualdade no tipo global, converta uma testemunha de
$\mathcal B_3(K)$ em átomos e aplique a versão global do
teorema~\ref{normalizacao:thm:norm}. A função global $G$ não muda. Na outra
direção, esqueça $\Agree$ e converta os átomos para células da
Fase 3 pelo lema~\ref{normalizacao:lem:witnesses}. A continuidade em $K$ é
preservada nas duas direções.
\end{proof}

\section{Orçamentos e erro com codomínio completo}
\label{normalizacao:sec:error}

% MAPA normalizacao:6.1
\begin{definicao}[orçamento da testemunha]\label{map:normalizacao:6.1}
Para $N,H:\N$, defina
\[
 \operatorname{Budget}_{N,H}(\Lambda):\Longleftrightarrow
 |\Lambda|\le N\ \land\ \forall a\in\Lambda,\ h(a)\le H.
\]
Com $H\ge2$, defina os subconjuntos de $\CK$
\begin{align*}
 \Aloc_{N,H}(K)&=\{g:\CK:\exists\Lambda,\
       \Loc_K(g,\Lambda)\land\operatorname{Budget}_{N,H}(\Lambda)\},\\
 \Astr_{N,H}(K)&=\{g:\CK:\exists\Lambda,\
       \Str_K(g,\Lambda)\land\operatorname{Budget}_{N,H}(\Lambda)\}.
\end{align*}
Os dois limites pertencem à mesma testemunha finita.
\end{definicao}

% MAPA normalizacao:6.2
\begin{teorema}[igualdade sob os mesmos recursos]\label{map:normalizacao:6.2}\label{normalizacao:thm:budget}
Para todos $N,H:\N$ com $H\ge2$,
$\Aloc_{N,H}(K)=\Astr_{N,H}(K)$.
\end{teorema}
\begin{proof}
A normalização produz um subconjunto da testemunha original, logo
preserva o limite de cardinalidade e cada cota $h(a)\le H$.
Isso dá a inclusão local para consolidada. A inclusão inversa usa
a própria testemunha consolidada e esquece $\Agree$.
\end{proof}

% MAPA normalizacao:6.3
\begin{definicao}[distância uniforme e ínfimo]\label{map:normalizacao:6.3}
Para $f,g:\CK$, defina o número real
\[
 d_K(f,g)=\max_{z:K}|f(z)-g(z)|\in\R_{\ge0}.
\]
O máximo existe porque a função sob o máximo é contínua e $K$ é
compacto não vazio. Ele é a norma uniforme de $f-g$.
Defina o reticulado completo
$\ERealPos=[0,+\infty]$, correspondente a \code{ENNReal}, e a inclusão
$\operatorname{ofReal}:\R\to\ERealPos$, que envia $x$ ao
valor finito $\max\{0,x\}$. Para argumentos não negativos, ela
coincide com a inclusão canônica.
Para $\mathcal S\subseteq\CK$, defina
\[
 \operatorname{Err}_K(f,\mathcal S)=
 \inf\{\operatorname{ofReal}(d_K(f,g)):g\in\mathcal S\}
 \quad\text{em }\ERealPos.
\]
Nesse reticulado, $\inf\varnothing=+\infty$. Ponha
\[
 E^{\rm loc}_{N,H}(f)=\operatorname{Err}_K(f,\Aloc_{N,H}(K)),\qquad
 E^{\rm con}_{N,H}(f)=\operatorname{Err}_K(f,\Astr_{N,H}(K)).
\]
\end{definicao}

% MAPA normalizacao:6.4
\begin{corolario}[igualdade dos erros e de suas condições de altura]\label{map:normalizacao:6.4}
Para todo $f:\CK$ e todos $N,H:\N$ com $H\ge2$,
\[
 E^{\rm loc}_{N,H}(f)=E^{\rm con}_{N,H}(f).
\]
Para $\varepsilon\ge0$, os conjuntos de alturas admissíveis
$\{H\ge2:E^{\rm loc}_{N,H}(f)\le\operatorname{ofReal}(\varepsilon)\}$
e seu análogo consolidado coincidem. Se são não vazios, seus mínimos
naturais coincidem.
\end{corolario}
\begin{proof}
O teorema~\ref{normalizacao:thm:budget} identifica exatamente os conjuntos sobre
os quais os ínfimos são tomados. A igualdade dos conjuntos de alturas
segue substituindo esses valores iguais; um subconjunto não vazio
de $\N$ possui um único menor elemento.
\end{proof}

% MAPA normalizacao:6.5
\begin{proposicao}[testemunhas já ativas]\label{map:normalizacao:6.5}\label{normalizacao:prop:fixed}
Se $\Act_K(a)$ para todo $a\in\Lambda$, então
$\Norm_K(\Lambda)=\Lambda$. Em particular, isso vale se
$J_a\subseteq K$ para todo $a\in\Lambda$.
\end{proposicao}
\begin{proof}
A primeira afirmação é a regra de pertinência do filtro. Para a
segunda, $\alpha\le L_a<X_a<R_a\le\beta$, logo
$X_a\in O_a\cap K^\circ$.
\end{proof}

Assim, nas testemunhas cujas células já estão dentro de $K$, a
normalização mantém cada rótulo e a própria família, além das
cotas de recursos. Para as demais testemunhas, a igualdade de classes
é obtida removendo somente átomos inativos.

\section{Ordenação finita e cadeia normalizada}
\label{normalizacao:sec:chain}

% MAPA normalizacao:7.1
\begin{definicao}[ordem auxiliar total dos átomos]\label{map:normalizacao:7.1}\label{normalizacao:def:sort}
Compare átomos pela ordem lexicográfica da chave
\[
 \zeta(a)=\bigl(L_\Q(C(a)),a.k,a.m,a.y,a.s\bigr)
          \in\Q\times\N\times\Z\times\Q\times\Q.
\]
Essa chave é injetiva, pois contém os quatro campos de dados.
Ela induz uma ordem linear $\triangleleft$ em $\Atoms$.
Defina $\operatorname{sort}:\FS(\Atoms)\to\List(\Atoms)$ por
recursão bem fundada na cardinalidade:
\[
 \operatorname{sort}(\varnothing)=[],\qquad
 \operatorname{sort}(F)=a::\operatorname{sort}(F\setminus\{a\})
       \quad(F\ne\varnothing),
 \qquad a=\min_{\triangleleft}F.
\]
O mínimo existe e é único porque $F$ é finito não vazio e a ordem
é linear. Como $a\in F$,
$|F\setminus\{a\}|=|F|-1$, o que fornece a medida de término.

Para listas, $[]$ é a lista vazia e $a::L$ acrescenta $a$ à frente
de $L$. Comprimento e acesso satisfazem
\[
 \operatorname{length}([])=0,\quad
 \operatorname{length}(a::L)=1+\operatorname{length}(L),\quad
 (a::L)_0=a,\quad (a::L)_{j+1}=L_j,
\]
sempre sob o limite de índice correspondente.
\end{definicao}

% MAPA normalizacao:7.2
\begin{lema}[propriedades da lista ordenada]\label{map:normalizacao:7.2}\label{normalizacao:lem:sort}
$\operatorname{sort}(F)$ tem comprimento $|F|$, não tem repetições,
contém exatamente os elementos de $F$ e é estritamente crescente
em $\triangleleft$. Se $\Sep(F)$, seus extremos esquerdos são
estritamente crescentes.
\end{lema}
\begin{proof}
As quatro primeiras propriedades seguem por indução em $|F|$:
o caso vazio é imediato; no passo, o mínimo pertence a $F$, é
removido antes da chamada recursiva e precede estritamente todos
os elementos restantes. A fórmula do comprimento dá $|F|$.

Sob $\Sep(F)$, dois átomos distintos não podem ter o mesmo
extremo esquerdo. Se $L_a=L_b=L$, o intervalo
$(L,\min\{R_a,R_b\})$ é não vazio e está em $O_a\cap O_b$.
Como a primeira coordenada de $\zeta$ é $L_\Q$, a ordem da lista
é então estrita também nessa coordenada.
\end{proof}

% MAPA normalizacao:7.3
\begin{definicao}[tipo dos dados de cadeia]\label{map:normalizacao:7.3}\label{normalizacao:def:chain}
Para $\Lambda:\FS(\Atoms)$, uma estrutura
$\operatorname{NormChain}_K(\Lambda)$ contém os dados
\[
 M:\N,\quad h_M:1\le M,\quad
 a:\Fin(M)\to\Atoms,\quad t:\Fin(M+1)\to\Q
\]
e as seguintes propriedades:
\begin{enumerate}
\item $a$ é injetiva e
$\operatorname{image}(a,\operatorname{univ}_M)=\Norm_K(\Lambda)$;
\item $t_j<t_{j+1}$ para $0\le j<M$;
\item $L_\Q(C(a_j))=t_j$ e $R_\Q(C(a_j))=t_{j+1}$ para $j<M$;
\item $\uparrow t_0\le\alpha$ e $\beta\le\uparrow t_M$;
\item $\alpha<\uparrow t_j<\beta$ para $1\le j<M$.
\end{enumerate}
Na segunda e terceira cláusulas, o índice $j$ de $t$ é a inclusão
$\langle j,h_j\rangle\mapsto\langle j,h_j':j<M+1\rangle$;
o índice seguinte é $\langle j+1,j+1<M+1\rangle$.
Por definição dos extremos, a terceira cláusula implica
$J_{a_j}=[\uparrow t_j,\uparrow t_{j+1}]$, com ambos os extremos incluídos.
\end{definicao}

% MAPA normalizacao:7.4
\begin{teorema}[construção determinada da cadeia]\label{map:normalizacao:7.4}\label{normalizacao:thm:chain}
De $\Loc_K(g,\Lambda)$ obtemos uma estrutura
$\operatorname{NormChain}_K(\Lambda)$ pela seguinte construção:
\begin{align*}
 F&:=\Norm_K(\Lambda),& M&:=|F|,\\
 a_j&:=(\operatorname{sort}(F))_j\quad(j<M),\\
 t_0&:=L_\Q(C(a_0)),&
 t_{j+1}&:=R_\Q(C(a_j))\quad(j<M).
\end{align*}
Não há escolhas de enumeração adicionais nessa construção.
\end{teorema}
\begin{proof}
O teorema~\ref{normalizacao:thm:norm} dá $F\ne\varnothing$, logo $M\ge1$
e o índice zero de $a$ existe. O lema~\ref{normalizacao:lem:sort} garante
que $a$ enumera $F$ bijetivamente e que $L_{a_j}<L_{a_{j+1}}$.
Para quaisquer $i<j$, o lema~\ref{normalizacao:lem:separated} dá
$R_{a_i}\le L_{a_j}$: a outra ordenação contradiz a ordem
estrita dos extremos esquerdos.

Para índices consecutivos, suponha
$R_{a_j}<L_{a_{j+1}}$. Seja
$z=(R_{a_j}+L_{a_{j+1}})/2$. Escolha pontos ativos
$u\in O_{a_j}\cap K^\circ$ e
$v\in O_{a_{j+1}}\cap K^\circ$. Temos
$\alpha<u<z<v<\beta$, logo $z\in K$.
Para $i\le j$, a ordenação fornece $R_{a_i}\le R_{a_j}<z$;
para $i\ge j+1$, temos $z<L_{a_{j+1}}\le L_{a_i}$.
Assim nenhuma célula de $F$ contém $z$, contrariando sua cobertura.
Concluímos
\[
 R_{a_j}=L_{a_{j+1}}\qquad(0\le j<M-1).
\]
Isso prova a terceira cláusula da cadeia e, como cada largura é
positiva, também a segunda. A cobertura de $\alpha$ e de $\beta$
fornece $L_{a_0}\le\alpha$ e $\beta\le R_{a_{M-1}}$.
Cada nó interno é a interseção de duas células distintas de $F$;
o teorema~\ref{normalizacao:thm:norm} o situa em $K^\circ$.
Todos os nós são racionais por sua fórmula e diádicos pelo
lema~\ref{normalizacao:lem:geometry}. Cada campo de
$\operatorname{NormChain}$ foi assim construído e justificado.
\end{proof}

\section{Extensão ReLU por recursão finita}
\label{normalizacao:sec:relu}

% MAPA normalizacao:8.1
\begin{definicao}[parâmetros e funções parciais]\label{map:normalizacao:8.1}\label{normalizacao:def:relu}
Dada uma cadeia $q:\operatorname{NormChain}_K(\Lambda)$, use seus
campos $M,a,t$. Defina
\[
 s_j=a_j.s\in\Q\quad(j<M),\qquad
 c_0=u(a_0)\in\Q,\qquad
 w_j=s_j-s_{j-1}\in\Q\quad(1\le j<M),
\]
\[
 \rho:\R\to\R,\qquad\rho(x)=\max\{0,x\}.
\]
Para cada natural $r<M$, defina $F_r:\R\to\R$ por recursão em $r$:
\begin{align*}
 F_0(x)&=\uparrow c_0+\uparrow s_0\,x,\\
 F_{r+1}(x)&=F_r(x)+\uparrow(s_{r+1}-s_r)\,
                       \rho(x-\uparrow t_{r+1})
                    \quad(r+1<M).
\end{align*}
A segunda cláusula só usa índices menores que $M$. Em uma definição
Lean dependente, o argumento é $(r:\N)$ junto à prova $r<M$;
na chamada recursiva $r+1\mapsto r$, essa prova é enfraquecida
para $r<M$. A medida de término é $r$.
Defina
\[
 R(q):=F_{M-1}:\R\to\R.
\]
O índice final é válido porque $M\ge1$.
\end{definicao}

% MAPA normalizacao:8.2
\begin{lema}[expansão da recursão]\label{map:normalizacao:8.2}\label{normalizacao:lem:sum}
Para $r<M$,
\[
 F_r(x)=\uparrow c_0+\uparrow s_0\,x+
        \sum_{1\le j\le r}\uparrow w_j\,\rho(x-\uparrow t_j).
\]
Em particular,
\begin{align}
 R(q)(x)&=\uparrow c_0+\uparrow s_0\,x+
        \sum_{1\le j<M}\uparrow w_j\,\rho(x-\uparrow t_j)
                                                       \label{normalizacao:eq:relu_sum}\\
 &=\ell_{a_0}(\uparrow t_0)+\uparrow s_0(x-\uparrow t_0)+
        \sum_{1\le j<M}\uparrow w_j\,\rho(x-\uparrow t_j).
                                                       \label{normalizacao:eq:relu_anchor}
\end{align}
Se $M=1$, $R(q)=\ell_{a_0}$ em toda a reta.
\end{lema}
\begin{proof}
A primeira identidade é indução em $r$: a soma inicial é vazia,
e a cláusula sucessora acrescenta exatamente o termo $j=r+1$.
Tome $r=M-1$ para obter \eqref{normalizacao:eq:relu_sum}. Pelo
lema~\ref{normalizacao:lem:line},
$\ell_{a_0}(\uparrow t_0)=\uparrow c_0+\uparrow s_0\,\uparrow t_0$,
que dá \eqref{normalizacao:eq:relu_anchor}. Para $M=1$, não existe termo de
correção e a função inicial é a primeira reta.
\end{proof}

% MAPA normalizacao:8.3
\begin{lema}[identidade de duas retas incidentes]\label{map:normalizacao:8.3}\label{normalizacao:lem:incident}
Se $j+1<M$ e
$\ell_{a_{j+1}}(\uparrow t_{j+1})=
 \ell_{a_j}(\uparrow t_{j+1})$, então, para todo $x:\R$,
\[
 \ell_{a_{j+1}}(x)-\ell_{a_j}(x)
   =\uparrow(s_{j+1}-s_j)(x-\uparrow t_{j+1}).
\]
\end{lema}
\begin{proof}
Escreva as duas retas como intercepto mais inclinação vezes $x$.
A igualdade em $\uparrow t_{j+1}$ determina a diferença dos
interceptos como $-\uparrow(s_{j+1}-s_j)\uparrow t_{j+1}$.
Substituir essa diferença e fatorar prova a identidade.
\end{proof}

% MAPA normalizacao:8.4
\begin{teorema}[correção da extensão]\label{map:normalizacao:8.4}\label{normalizacao:thm:extension}
Suponha $\Loc_K(g,\Lambda)$ e seja
$q:\operatorname{NormChain}_K(\Lambda)$ qualquer cadeia válida.
Então $R(q)$ é contínua em $\R$ e satisfaz
\begin{align*}
 R(q)(x)&=\ell_{a_j}(x)
            &&(j<M,\ \uparrow t_j\le x\le\uparrow t_{j+1}),\\
 R(q)(\iota_K(z))&=g(z)&& (z:K),\\
 R(q)(x)&=\ell_{a_0}(x)&& (x\le\uparrow t_0),\\
 R(q)(x)&=\ell_{a_{M-1}}(x)&& (x\ge\uparrow t_{M-1}).
\end{align*}
Todos os parâmetros de \eqref{normalizacao:eq:relu_sum} são racionais e seus
nós internos pertencem a $K^\circ\cap\D$. Se $M>1$, a igualdade
com a primeira reta vale também para $x\le\uparrow t_1$.
\end{teorema}
\begin{proof}
O teorema~\ref{normalizacao:thm:norm} fornece $\Agree$ na família normalizada.
Como as duas células adjacentes contêm $\uparrow t_{j+1}$,
o lema~\ref{normalizacao:lem:incident} aplica-se a cada par consecutivo.

Por indução, para $r<M$,
\[
 F_r(x)=\ell_{a_r}(x)\quad\text{se }x\ge\uparrow t_r.
\]
O caso $r=0$ decorre da forma global da primeira reta.
No passo, $x\ge\uparrow t_{r+1}>\uparrow t_r$ permite usar a
hipótese de indução, e
$\rho(x-\uparrow t_{r+1})=x-\uparrow t_{r+1}$. A identidade
de retas incidentes produz precisamente $\ell_{a_{r+1}}(x)$.

Se $x\in[\uparrow t_j,\uparrow t_{j+1}]$, a identidade anterior
vale para $F_j(x)$. Todos os termos posteriores têm nós
$\uparrow t_r\ge\uparrow t_{j+1}$ e, portanto, são zero nesse
ponto. Segue a primeira conclusão. A cobertura por essas células
e $\Match_K(g,\Norm_K(\Lambda))$ dão a segunda.
À esquerda de $\uparrow t_0$, todas as correções são zero; quando
$M>1$, o mesmo vale até $\uparrow t_1$. À direita de
$\uparrow t_{M-1}$, a identidade indutiva para $F_{M-1}$ dá a
última reta. O caso $M=1$ é o lema~\ref{normalizacao:lem:sum}.

Cada $F_r$ é contínua por indução: a inicial é afim, $\rho$ é
contínua como máximo de funções contínuas, e soma, translação e
produto por constante preservam continuidade. A racionalidade
vem dos tipos de $c_0,s_0,w_j,t_j$. A localização e a diadicidade
dos nós foram provadas no teorema~\ref{normalizacao:thm:chain}.
\end{proof}

% MAPA normalizacao:8.5
\begin{lema}[unicidade de uma reta a partir de dois pontos]\label{map:normalizacao:8.5}\label{normalizacao:lem:two_points}
Se duas funções afins coincidem em dois pontos distintos, coincidem
em toda a reta.
\end{lema}
\begin{proof}
Subtraia as duas igualdades de valores. O resultado é
$(s-s')(x_2-x_1)=0$. Como $x_1\ne x_2$, as inclinações são
iguais; qualquer uma das igualdades originais dá a igualdade dos
interceptos.
\end{proof}

% MAPA normalizacao:8.6
\begin{teorema}[independência da testemunha]\label{map:normalizacao:8.6}\label{normalizacao:thm:unique}
Se $\Loc_K(g,\Lambda)$ e $\Loc_K(g,\Lambda')$, e se
$q:\operatorname{NormChain}_K(\Lambda)$ e
$q':\operatorname{NormChain}_K(\Lambda')$, então
$R(q)=R(q')$ como funções $\R\to\R$.
\end{teorema}
\begin{proof}
As extensões coincidem em $K$ pelo teorema~\ref{normalizacao:thm:extension}.
Para cada cadeia, defina o número positivo
\[
 \eta_L(q)=\begin{cases}
 \uparrow t_1-\alpha,&M>1,\\
 \beta-\alpha,&M=1.
 \end{cases}
\]
A positividade decorre da localização dos nós internos e de $h_K$.
Ponha
$\delta_L=\tfrac12\min\{\beta-\alpha,\eta_L(q),\eta_L(q')\}>0$.
As duas extensões são suas respectivas primeiras retas em
$(-\infty,\alpha+\delta_L]$. Elas coincidem em
$x_1=\alpha+\delta_L/3$ e
$x_2=\alpha+2\delta_L/3$, dois pontos distintos de $K$.
O lema~\ref{normalizacao:lem:two_points} identifica as primeiras retas
em toda a reta, logo as extensões coincidem para $x\le\alpha$.

À direita, use
$\eta_R(q)=\beta-\uparrow t_{M-1}$ se $M>1$, e
$\eta_R(q)=\beta-\alpha$ se $M=1$; defina $\delta_R$
pelo mesmo mínimo e use os pontos
$\beta-\delta_R/3$ e $\beta-2\delta_R/3$.
As últimas retas coincidem, logo as extensões coincidem para
$x\ge\beta$. Esses dois conjuntos e $K$ cobrem $\R$.
\end{proof}

% MAPA normalizacao:8.7
\begin{definicao}[operador de extensão das montagens]\label{map:normalizacao:8.7}
Para o subtipo $\{g:\CK\mid g\in\Aloc(K)\}$, defina
\[
 \mathcal E_K:\{g:\CK\mid g\in\Aloc(K)\}\longrightarrow C(\R,\R)
\]
como a única extensão $R(q)$ produzida por uma representação local
e uma cadeia válida. A existência usa os
teoremas~\ref{normalizacao:thm:chain} e \ref{normalizacao:thm:extension}; a unicidade é o
teorema~\ref{normalizacao:thm:unique}. Logo a definição independe da prova de
pertinência, da testemunha e da cadeia escolhidas. Vale
$\Res_K(\mathcal E_K(g))=g$.
\end{definicao}

Em \eqref{normalizacao:eq:relu_anchor}, o valor inicial é
$\ell_{a_0}(\uparrow t_0)$, definido em toda a reta. Também é
$\mathcal E_K(g)(\uparrow t_0)$. A avaliação direta de $g(t_0)$
só é uma expressão bem tipada quando $\uparrow t_0\in K$.

\clearpage
\setcounter{section}{0}\renewcommand{\thesection}{II.\Alph{section}}
\section{Vinculação às definições Lean existentes}
\label{normalizacao:sec:adapters}

A correspondência abaixo fixa o significado dos nomes do pacote
\code{Biblioteca_Lean_Equivalencias_v2_COMPLETO.zip}. Ela é uma
especificação de integração: os lemas de identificação devem ser
provados ao reutilizar as definições existentes. A definição matemática
não fica a cargo de uma nova escolha do formalizador.

\begin{longtable}{@{}>{\raggedright\arraybackslash}p{0.28\linewidth}>{\raggedright\arraybackslash}p{0.67\linewidth}@{}}
\toprule
Objeto neste texto & Identificador ou expressão no pacote examinado\\
\midrule
\endhead
$\Atoms$ e seus campos & \code{Dyadic.Lib.Atom}; os campos $y,s$ são
\code{y0,s}, e a prova de não nulidade é \code{hs}.\\
$L_\Q,R_\Q,X_\Q$ do átomo &
\code{Dyadic.Lib.Atom.left}, \code{right}, \code{center}.\\
$J_a$, $\ell_a$, $u(a)$, $h(a)$ &
\code{Dyadic.Lib.Atom.cell}, \code{line}, \code{intercept}, \code{height}.\\
Célula das Fases 2 e 3 &
\code{Dyadic2.cellJ k m}; \code{Fase3.cell k m}.\\
Centro das Fases 2 e 3 &
\code{Dyadic2.center k m}; \code{Fase3.center k m}.\\
$\rankq_2,\rankq_2^*$ &
\code{Dyadic2.Qrank}; \code{Dyadic2.Qstarrank}.\\
$Q_2,Q_2^*$ & \code{Dyadic2.Qenum}; \code{Dyadic2.Qstar}.\\
$\rankq_3,\rankq_3^*$ &
\code{Fase3.rankQ}; \code{Fase3.rankQstar}.\\
$Q_3,Q_3^*$ & \code{Fase3.Qenum}; \code{Fase3.QstarEnum}.\\
$\rankq_{\rm c},\rankq_{\rm c}^*$ &
\code{Dyadic.Lib.rankQ}; \code{Dyadic.Lib.rankQstar}.\\
$\enc_{\rm c},\dec_{\rm c}$ &
\code{Dyadic.Lib.encode}; \code{Dyadic.Lib.decode}.\\
Átomo decodificado da Fase 2 & \code{SemanticBridge.ofFase2 n}; campos
\code{Dyadic2.scale n}, \code{cellIdx n}, \code{y0n n}, \code{sn n}.\\
Átomo decodificado da Fase 3 & \code{SemanticBridge.ofFase3 n}; campos
\code{Fase3.gK n}, \code{gM n}, \code{y0 n}, \code{slope n}.\\
$\lambda_2(n,\cdot)$, $\lambda_3(n,\cdot)$ &
\code{Dyadic2.line n}; \code{Fase3.atomLine n}.\\
$\Loc^G_K(G,\Lambda)$ &
\code{SemanticBridge.AssemblyCore K G Lambda}.\\
Local com concordância interna redundante &
\code{SemanticBridge.IsAssemblyRepOn K G Lambda}; equivale a
\code{AssemblyCore} por \code{SemanticBridge.local_iff_core}.\\
$\Str^G_K(G,\Lambda)$ &
\code{Dyadic.Lib.IsAssemblyRep K G Lambda}.\\
$\mathcal B_2(K)$ & \code{Dyadic2.Asm alpha beta}.\\
$\mathcal B_3(K)$ & Predicado \code{Fase3.IsAssembly alpha beta G}.\\
$\mathcal B_{\rm con}(K)$ & \code{Dyadic.Lib.Asm (Set.Icc alpha beta)}.\\
Classe consolidada com orçamentos &
\code{Dyadic.Lib.AsmNH (Set.Icc alpha beta) N H}, no tipo global;
a versão em $\CK$ é sua imagem por restrição.\\
Condição de interseções internas &
\code{SemanticBridge.IntersectionsInside K Lambda}; a terceira conclusão
do teorema~\ref{normalizacao:thm:norm} implica esse predicado.\\
Conversão local para consolidada sob esse critério &
\code{SemanticBridge.strong_iff_local_of_intersections_inside}.\\
\bottomrule
\end{longtable}

\paragraph{Predicados auxiliares importados.}
Para tornar explícitas as duas linhas auxiliares da tabela, defina
\begin{align*}
 \operatorname{InnerAgree}_K(\Lambda)
 &:\Longleftrightarrow
 \forall a,b\in\Lambda\ \forall x\in J_a\cap J_b\cap K,\
                    \ell_a(x)=\ell_b(x),\\
 \operatorname{IntersectionsInside}_K(\Lambda)
 &:\Longleftrightarrow
 \forall a,b\in\Lambda,\ a\ne b\Rightarrow J_a\cap J_b\subseteq K.
\end{align*}
O predicado importado \code{IsAssemblyRepOn K G Lambda} é
$\Loc^G_K(G,\Lambda)\land\operatorname{InnerAgree}_K(\Lambda)$,
com a associação de conjunções expandida na sua definição original.
A concordância interna deriva de $\Match^G$. Sob
$\operatorname{IntersectionsInside}_K(\Lambda)$, ela implica
$\Agree(\Lambda)$, pois encontros de átomos iguais têm igualdade
reflexiva e encontros de átomos distintos estão em $K$.

\paragraph{Identidades obrigatórias dos adaptadores.}
Se o tipo existente de átomos é importado, $\Atoms$ deve ser um
alias desse tipo, com as projeções indicadas na tabela.
As duas construções \code{ofFase2} e \code{ofFase3} são fixadas
por seus quatro campos de dados. As identidades de integração
requeridas são
\[
 \begin{aligned}
 \text{\ttfamily SemanticBridge.ofFase2 n}&=\dec_2(n),\\
 \text{\ttfamily SemanticBridge.ofFase3 n}&=\dec_3(n).
 \end{aligned}
\]
Para $i=2,3$, devem ser transportadas também as identidades
\[
 J_{\dec_i(n)}=J(C_i(n)),\qquad
 \ell_{\dec_i(n)}=\lambda_i(n,\cdot).
\]
As primeiras duas igualdades seguem expandindo pareamento, enumeração dos
inteiros e inversas dos postos definidos na
seção~\ref{normalizacao:sec:enumeration}. As identidades de suporte e reta seguem das
seções~\ref{normalizacao:sec:objects} e \ref{normalizacao:sec:enumeration}.
A equivalência dos predicados é então provada campo a campo,
incluindo as coerções entre $x:\R$ com prova $x\in K$ e $z:K$.

O enunciado final no tipo global é
\[
 \forall G:\R\to\R,\qquad
 \text{\ttfamily Fase3.IsAssembly}\ \alpha\ \beta\ G
 \ \Longleftrightarrow\
 G\in\text{\ttfamily Dyadic.Lib.Asm}\bigl([\alpha,\beta]\bigr),
\]
sob $\alpha<\beta$. A igualdade com \code{Dyadic2.Asm} é
expressa no tipo $\CK$ por $\ImageRes_K$, como no
teorema~\ref{normalizacao:thm:bridge}. Os dois tipos de função permanecem
explícitos nos enunciados.

\section{Ordem de formalização e fronteiras das obrigações}
\label{normalizacao:sec:dependencies}

\begin{longtable}{@{}>{\raggedright\arraybackslash}p{0.07\linewidth}>{\raggedright\arraybackslash}p{0.28\linewidth}>{\raggedright\arraybackslash}p{0.59\linewidth}@{}}
\toprule
Etapa & Dados definidos & Propriedades necessárias antes da etapa seguinte\\
\midrule
\endhead
1 & Subtipos, $K$, $\CK$, conjuntos finitos &
Extensionalidade; $\alpha<\beta$; compacidade e
$K=\cls(K^\circ)$.\\
2 & Racionais reduzidos, $\Cell$, $\Atoms$, $h$ &
Largura positiva; interior não vazio; igualdade de células e de retas.\\
3 & $\Sep,\Cov$,\newline $\Match,\Agree$,\newline $\Loc,\Str$ &
Cada campo tem domínio e quantificadores fixados;
concordância interna deriva de $\Match$.\\
4 & $\Act$, filtro $\Norm$ &
Cobertura dos extremos; interseções internas; contrato forte;
idempotência e não aumento de cardinalidade.\\
5 & $T,\pi,Z$, postos, $Q_i,Q_i^*,\enc_i,\dec_i$ &
Inversas nos dois sentidos; suporte e reta preservados.\\
6 & $\operatorname{CellRep}$, $T_i,I,\sigma,e_{i,\Lambda}$,
$\Res,E_0$ &
Conversões válidas; cardinalidades; restrição bem tipada;
identificação com as definições importadas.\\
7 & Classes limitadas, $d_K$, $\operatorname{Err}_K$ &
Igualdade das classes com os mesmos $N,H$;
ínfimo em \code{ENNReal}, inclusive na classe vazia.\\
8 & Ordem $\zeta$, lista \code{sort}, $\operatorname{NormChain}$ &
Término por cardinalidade; bijeção finita; adjacência sem lacunas;
limites de todos os índices.\\
9 & $F_r$, $R(q)$, $\mathcal E_K$ &
Término por $r$; identidade de retas incidentes; continuidade;
coincidência em $K$ e independência da testemunha.\\
\bottomrule
\end{longtable}

A normalização e as igualdades das classes encerram-se nas etapas 4 e 6.
A extensão ReLU depende dessas etapas e pode ser formalizada depois,
em lemas separados.

As inversas únicas e as decisões clássicas sobre reais podem ser definidas
com \code{noncomputable}. A marca não comprova seus contratos:
decodificação, suporte, cobertura e igualdade de funções continuam
exigindo provas. As equivalências não exigem computabilidade.


